\begin{lemma}\label{lem:octic-315-exact-specialization}
Let $F=\mathbb Q(\zeta_3)$, $X=X_1(3,15)$, and
$J=\operatorname{Jac}(X)$. Fix an $F$-rational cusp $c_0$, and
let $C\subseteq J(F)$ be generated by $F$-rational cuspidal
divisors of degree zero. Reduction at the prime above $2$ is
injective on $J(F)$. Every $F$-rational noncuspidal effective
divisor $D$ of degree $4$ satisfies
\[
 [D-4c_0]\in C.
\]
There are at most twelve possible classes of such divisors.
The twelve classes form one orbit under the $F$-defined
automorphisms of $X$, with the action on degree-$4$ divisor
classes referenced to $4c_0$.
\end{lemma}
\begin{proof}
Lemma~\ref{lem:octic-315-cusp-image} gives
\[
 C\simeq[3,3,24,96,96],\qquad\#C=1990656,
 \qquad\#(J(F)/C)\mid3.
\]
Here $J(F)$ is finite by the rank-zero theorem of
\cite[Theorem~4.1]{DS}. The cuspidal image at $\mathbb F_4$
has the same order as $C$, so reduction is injective on $C$.
Moreover $J(F)[2^\infty]=C[2^\infty]$, since the quotient has
odd order. Prime-to-$2$ torsion injects under good reduction,
so reduction is injective on all of $J(F)$.

For completeness, the index bound uses the exact Jacobian
orders at the split primes $7,13,19$:
\[
149299200,\quad12230590464,\quad247669456896.
\]
Their gcd is $5971968=3\#C$. Good unramified odd-prime
reduction is injective on torsion. All cusps split over
$\mathbb Q(\zeta_{15})$, and the Frobenius at $7$ generates
its degree-$4$ Galois group over $F$. Thus the twenty cusp
blocks of degrees $1^{16},4^4$ at $7$ are precisely the
$F$-rational cusp orbits. The integral-unit signatures in
Lemma~\ref{lem:octic-315-cusp-image} match these same global
divisors at $\mathbb F_4$ and $\mathbb F_7$. The common
free cusp-divisor maps have equal kernels and image of
order $1990656$ at both primes.

We enumerate every effective degree-$4$ divisor at
$\mathbb F_4$: $5000$ entirely cuspidal divisors and $102$
noncuspidal quartic places. At $7$ the entirely cuspidal
branch is impossible for a global noncuspidal $D$ by the
odd-prime cuspidal-lifting argument, rank zero, and
$\operatorname{gon}_F(X)\geq6>4$. The remaining patterns
are a quadratic place with a cuspidal degree-$2$ complement,
a cubic place with a rational cusp, a quartic place, and
two quadratic places, including repetition. Their numbers
are respectively
\[
24\binom{17}{2}=3264,\quad144\cdot16=2304,
\quad396,\quad\binom{25}{2}=300.
\]
All complements occur. In particular this enumeration uses
all $5102$ and $6264$ local divisors, rather than selected
representatives.

Since $3J(F)\subseteq C$, the two reductions of
$3[D-4c_0]$ must belong to the diagonal image of these same
global cusp divisors. The exact enumeration leaves sixty
pairs of local divisor classes. All sixty $\mathbb F_4$
classes lie in the image of $C$. Injectivity of reduction
therefore implies $[D-4c_0]\in C$: subtract the unique
matching global cuspidal class and use injectivity on
$J(F)$.

We may now compare the unmultiplied classes in the diagonal
cuspidal subgroup. Exactly twelve pairs remain, giving
twelve distinct elements of $C$. The file
\texttt{M3N15\_D4\_exact\_cusp\_lifts.txt} records an explicit
signed global cuspidal representative
\[
 L=4c_0+\sum_{i=1}^{20}v_i(B_i-\deg(B_i)c_0),
 \qquad v_i\in\{0,1,-1\},
\]
for each class. Thus the assertion does not require full
cuspidal generation of $J(F)$.

Finally put $Q_3=5Q$ and $Q_5=3Q$, so
$Q=2Q_3+2Q_5$. The $48$ automorphisms
\[
 (P,Q)\longmapsto
 \bigl(aP+bQ_3,\ 2(cP+dQ_3)+2uQ_5\bigr),
 \quad
 \begin{pmatrix}a&b\\c&d\end{pmatrix}\in
 \operatorname{SL}_2(\mathbb F_3),\quad u\in\{1,2\},
\]
after simultaneous negation, act transitively on these
twelve classes. The calculation applies all $48$ maps to
an actual quadratic-plus-quadratic divisor at $7$, and
compares the resulting classes with the twelve global cusp
lifts. The class action is affine:
\[
 a\longmapsto g(a)+4[g(c_0)-c_0],
\]
and this translation is included by computing each image
divisor minus $4c_0$. Odd-prime injectivity on $J(F)$ lifts
the equality of classes to $F$.

The independent integer-Hermite-normal-form replay
\texttt{replay\_exact\_3\_15\_D4.py} verifies the actual finite
group relations, the sixty cuspidal containments, the twelve
unmultiplied matches and signed lifts, and all $48$ class
images. Since $\operatorname{gon}_F(X)>4$, each of these
classes has at most one effective $F$-rational degree-$4$
representative. The lemma gives a finite classification
problem; it does not assert that any of the twelve classes
has such a representative.
\end{proof}
